% Reproduced inherited proof; finding P-F26.
\subsection{Finite lifts and the entire smooth fixed locus}

\begin{proposition}[P-F26: finite two-power lifts and the smooth chain locus]
\label{prop:finitecubiclift}
Let $n\ge4$, $M=2^{n-1}$, $g=\gcd(n,3)$, and
$X=\operatorname{Spec}T^{(n)}$ with the canonical stretch grading
of P-F19. Write $\Delta$ for the discriminant, viewed as a
homogeneous function on $X$, and let $p_n$ be the quotient point
of the chain cubic $F_n$. Then the following statements hold.

\begin{enumerate}
\item Every nonzero closed point $p\in X^{\boldsymbol\mu_M}$
has a representative cubic $F$ with closed $SL_n$ orbit and an
exact-form stabilizer element $A\in GL_n$ of finite two-power
order such that $\det A$ has order exactly $M$.
\item If $\Delta(p)\ne0$, the lift $A$ can be chosen of order
exactly $M$, and $F$ is $GL_n$-equivalent to $F_n$.
\item Scheme-theoretically, the part of the fixed scheme in
the smooth-hypersurface locus is the single grading orbit:
\[
 X^{\boldsymbol\mu_M}\cap D(\Delta)
      =\mathbb G_m\cdot p_n.
 \tag{P-F26a}
\]
This is a smooth reduced curve. Every irreducible component of
$X^{\boldsymbol\mu_M}$ other than the chain component of P-F24
is contained in $V(\Delta)$.
\end{enumerate}
The first statement does not assert that the lift has order $M$
when $F$ is singular.
\end{proposition}

\begin{proof}
Put $\gamma=3/g$, an odd integer. Recall the actual ring
identification from P-F19:
\[
 T^{(n)}_s=\bigl(\operatorname{Sym}^{(n/g)s}
                         (\operatorname{Sym}^3\mathbb C^n)\bigr)^{SL_n}.
\]
The $GL_n$-character on this subspace is $(\det)^{\gamma s}$.
On points, with the convention $(B\cdot F)(x)=F(B^{-1}x)$,
one consequently has
\[
 P(B\cdot F)=(\det B)^{-\gamma s}P(F)
       \quad(P\in T^{(n)}_s).
 \tag{P-F26b}
\]
Fix a primitive $M$th root $\zeta$. Since $\gamma$ is odd,
there is a unique $\eta\in\boldsymbol\mu_M$ with
$\eta^{-\gamma}=\zeta$, and $\eta$ is primitive. Choose
$B\in GL_n$ with determinant $\eta$. Thus the action of $B$
on the quotient is the grading action with parameter $\zeta$.
This choice obtains a determinant of order $M$ directly; a
scalar-only lift could instead introduce an unnecessary odd
factor into its determinant order.

Every fiber of the affine reductive quotient contains exactly
one closed orbit. A precise primary source for this fact is
Luna, \emph{Slices \'etales}, Section~I.1, printed p.~84.
Choose $F$ in that closed orbit over $p$. The form $B\cdot F$
has a closed $SL_n$ orbit, since $B$ normalizes $SL_n$.
It has the same quotient point by (P-F26b) and the fixed-point
hypothesis. Hence $B\cdot F=S\cdot F$ for some $S\in SL_n$.
Replacing $B$ by $S^{-1}B$ produces a matrix fixing $F$
exactly and still of determinant $\eta$. Denote this matrix
temporarily by $C$.

We next replace $C$ by a finite two-power element without
changing its determinant. Let $C=C_sC_u$ be its commuting
semisimple and unipotent parts. Their actions on cubic forms
are again semisimple and unipotent. Decompose $F$ into the
eigenspaces of $C_s$ in that representation. On an eigenspace
of eigenvalue $a\ne1$, the operator $aC_u-1$ is invertible,
because every eigenvalue of $C_u$ is one. Since $C F=F$,
the components in those eigenspaces are zero. It follows that
$C_sF=F$; also $\det C_s=\det C=\eta$.

Diagonalize $C_s$, with eigenvalues $\alpha_1,\ldots,\alpha_n$
on the underlying vector space. Let
$\Gamma=\langle\alpha_1,\ldots,\alpha_n\rangle\subset\mathbb C^*$.
This is a finitely generated abelian group. Its torsion
subgroup is finite cyclic, and a choice of a free complement
gives a group homomorphism
\[
 \rho:\Gamma\longrightarrow\Gamma_{\mathrm{tors},2}
\]
that is the identity on its two-primary torsion subgroup,
and kills its odd torsion and the chosen free complement.
Here $\eta=\prod_i\alpha_i$ belongs to that two-primary
torsion subgroup, so $\rho(\eta)=\eta$.

Every monomial $x_1^{e_1}\cdots x_n^{e_n}$ with nonzero
coefficient in $F$ satisfies $\prod_i\alpha_i^{e_i}=1$;
inverting all the eigenvalues, as in the action on forms,
gives the same equality. Applying $\rho$ preserves this
relation. Therefore
\[
 A=\operatorname{diag}(\rho(\alpha_1),\ldots,
                       \rho(\alpha_n))
 \tag{P-F26c}
\]
fixes $F$ exactly, has finite two-power order, and has
determinant $\eta$. This proves the first statement.

Now assume $\Delta(p)\ne0$. The representative $F$ defines
a smooth hypersurface. Write the order of $A$ as $2^\ell$.
Because its determinant has order $M$, one has $\ell\ge n-1$.
At least one eigenvalue of $A$ has order $2^\ell$.
We give a direct smoothness argument for the reverse bound
and its equality case.

Suppose the eigenvalue on variable $x_i$ has nontrivial
two-power order. The coefficient of $x_i^3$ in an invariant
cubic must be zero. Thus the coordinate point $[e_i]$ lies
on the hypersurface. Smoothness there implies the presence
of a monomial $x_i^2x_j$ with nonzero coefficient for some
$j\ne i$. Invariance forces
\[
 \alpha_j=\alpha_i^{-2},
 \tag{P-F26d}
\]
where in this paragraph $\alpha_i$ denotes the eigenvalue
of the finite-order matrix $A$ itself. Its order is half
the order of $\alpha_i$. Starting with an eigenvalue of
order $2^\ell$ and iterating gives eigenvalues of orders
\[
 2^\ell,2^{\ell-1},\ldots,2,1.
\]
They occupy $\ell+1$ distinct variables. Hence $\ell+1\le n$.
Consequently $\ell=n-1$, and the chosen lift has order
exactly $M$.

In the equality case these are all the variables. After
permutation their eigenvalues are
\[
 \xi,\xi^{-2},\ldots,\xi^{(-2)^{n-1}}=1,
 \qquad \operatorname{ord}(\xi)=M.
 \tag{P-F26e}
\]
For completeness, the invariant degree-three monomials for
this spectrum are exactly those of the chain. The nonzero
residue weights $1,-2,4,\ldots,(-2)^{n-2}$ have distinct
two-adic valuations. In a vanishing sum of three weights,
a least nonzero valuation cannot occur only once. If it
occurs three times, the sum is not zero, since three is odd.
Thus it occurs twice, and the third weight must be minus
twice that weight, which is precisely the next chain weight.
At the last nonzero weight the next is zero. The all-zero
case gives the final cube. Therefore
\[
 F=c_1x_1^2x_2+\cdots+c_{n-1}x_{n-1}^2x_n+c_nx_n^3.
 \tag{P-F26f}
\]
Every $c_i$ is nonzero. For $i<n$ this follows again from
smoothness at $[e_i]$; if $c_n=0$, the final coordinate
point is singular as well. Starting with $b_n^3c_n=1$
and successively solving $b_i^2b_{i+1}c_i=1$ gives nonzero
complex numbers $b_i$. The diagonal substitution
$x_i\mapsto b_ix_i$ turns (P-F26f) into $F_n$. This proves
the second statement, including the equality classification.

The $GL_n$-action on the quotient is precisely the grading
action by (P-F26b): the map $\mathbb C^*\to\mathbb C^*$,
$z\mapsto z^{-\gamma}$, is surjective. Thus every fixed
closed point in $D(\Delta)$ belongs to the grading orbit
of $p_n$. Conversely that entire orbit is fixed by
$\boldsymbol\mu_M$ and lies in $D(\Delta)$, by P-F24 and
discriminant nonvanishing.

P-F24 also shows that the fixed scheme is smooth of dimension
one at every point of this orbit. Its intersection with
$D(\Delta)$ consequently is reduced and has exactly the
closed points of that orbit. Both carry their reduced
locally closed scheme structures, so they agree, proving
(P-F26a). More explicitly, any nonreduced or additional
finite-type locus would have a closed complex point,
contradicting the preceding local smoothness and point
classification. If a different irreducible fixed component
met $D(\Delta)$, it would pass through a smooth point of
the fixed scheme on this orbit. That point has a unique
component, namely the chain component of P-F24. This is a
contradiction. Thus all other components lie in $V(\Delta)$.
\end{proof}

\paragraph{Known order bound, and the precise remaining lift issue.}
For smooth hypersurfaces the two-power order bound also follows
from Zheng, \emph{On abelian automorphism groups of hypersurfaces},
arXiv:\texttt{2004.09008v3}, Theorem~4.8, printed pp.~11--12.
An exact-form stabilizer of two-power order has the same
projective order, since its scalar kernel is contained in
$\boldsymbol\mu_3$.
At degree three its factors $|1-(-2)^a|$ are odd, and the only
two-power factors in the theorem are $2^{b-1}$ with $b\le n$.
The elementary coordinate-point argument above proves the needed
bound and makes its equality case explicit. Theorem~4.3 alone
states an abstract subgroup embedding into the diagonal group
of a simple hypersurface; it should not be cited by itself as
conjugacy of the original representation. No historical priority
is claimed for the order bound or its chain realization.

For a singular closed-orbit representative $F$, put
$H=\operatorname{Stab}_{GL_n}(F)$. Once the determinant image
contains $\boldsymbol\mu_M$, an exact-order-$M$ lift exists
if and only if the sequence of complex groups
\[
 1\longrightarrow H\cap SL_n\longrightarrow
   \det^{-1}(\boldsymbol\mu_M)\cap H
   \xrightarrow{\det}\boldsymbol\mu_M\longrightarrow1
 \tag{P-F26g}
\]
has a group-theoretic section. Indeed an order-$M$ element
mapping to a generator defines a section, and the image
of a generator under a section has order exactly $M$.
Such splitting is not automatic for abstract group extensions:
the squaring map from a cyclic group of order $2M$ onto a
cyclic group of order $M$ has no such section when $M$ is a
power of two. This abstract example is not asserted to occur
as the stabilizer extension of a cubic. The proof above
establishes splitting in the smooth case. At this step it
confines all possible competing fixed components to the
discriminant boundary without assuming a splitting there.
P-F30 subsequently excludes every nonzero singular competitor
at this extremal determinant order, even for larger-order lifts;
P-F31 then proves uniqueness and noncancellation. The abstract
extension warning explains why that additional work is needed.

The primary texts are
\url{https://www.numdam.org/item/10.24033/msmf.110.pdf}
and \url{https://arxiv.org/pdf/2004.09008}.
